\section{Cipher Caesar}

Cipher Caesar adalah cipher substitusi monoalfabetik paling sederhana \cite{schneier1996,trappe2006}. Setiap huruf digeser sejumlah posisi tetap dalam alfabet. Dengan shift $k$, huruf $A$ menjadi huruf ke-$(1+k)$ mod 26.

\begin{table}[!htbp]
\centering
\caption{Tabel Caesar cipher dengan shift 3 (A$\to$D, B$\to$E, \ldots)}
\label{tab:caesar}
\begin{tabular}{@{}*{14}{c}@{}}
\toprule
Plain & A & B & C & D & E & F & G & H & I & J & K & L & M \\
Cipher & D & E & F & G & H & I & J & K & L & M & N & O & P \\
\midrule
Plain & N & O & P & Q & R & S & T & U & V & W & X & Y & Z \\
Cipher & Q & R & S & T & U & V & W & X & Y & Z & A & B & C \\
\bottomrule
\end{tabular}
\end{table}

Contoh dengan shift 3: HELLO $\to$ KHOOR. Dekripsi dilakukan dengan shift kebalikan (26 - k).

\begin{lstlisting}[style=cstyle,caption={Caesar Cipher dalam C (ringkas)}]
char caesar(char c, int shift) {
    if (c >= 'A' && c <= 'Z')
        return 'A' + (c - 'A' + shift) % 26;
    if (c >= 'a' && c <= 'z')
        return 'a' + (c - 'a' + shift) % 26;
    return c;
}
\end{lstlisting}

\begin{lstlisting}[style=pascalstyle,caption={Caesar Cipher dalam Pascal (ringkas)}]
function caesar_encrypt(s: string; k: integer): string;
var i: integer; result: string;
begin
  result := '';
  for i := 1 to length(s) do
    if (s[i] >= 'A') and (s[i] <= 'Z') then
      result := result + chr(ord('A') + (ord(s[i]) - ord('A') + k) mod 26)
    else
      result := result + s[i];
  caesar_encrypt := result;
end;
\end{lstlisting}

Kelemahan: hanya 26 kemungkinan kunci; mudah diserang dengan brute force. Implementasi lengkap dalam MATLAB dan C++ tersedia di Bab 14.
